% =============================================================================
% ELEMENTOS PRÉ-TEXTUAIS COM AS LOGOS OFICIAIS CINESOFT E CINEDOOM
% Projeto: Template Editorial — Manual de Processos de Software (CineSoft)
% =============================================================================

% 1. CAPA OFICIAL (PÁGINA 1)
\capamanual[images/capa.png]

% 2. FICHA INSTITUCIONAL E EXPEDIENTE COM AS LOGOS OFICIAIS (PÁGINA 2)
\thispagestyle{empty}
\vspace*{2pt}
\noindent
\begin{minipage}{0.48\textwidth}
  \includegraphics[height=1.2cm]{logos/cinesoft-logo-horizontal-color.png}
\end{minipage}%
\hfill
\begin{minipage}{0.48\textwidth}
  \raggedleft
  \includegraphics[height=1.2cm]{logos/cinedoom-logo-horizontal-color.png}
\end{minipage}

\vspace{10pt}
\noindent
{\fontsize{13pt}{17pt}\fontExtraBold\color{BrandNavy}CINESOFT SOLUTIONS — ENGENHARIA DE SOFTWARE}\\[3pt]
{\fontsize{10pt}{13.5pt}\fontMedium\color{BrandCrimson}Manual Corporativo de Processos de Ciclo de Vida de Software (Plataforma CineDoom)}

\vspace{12pt}
\noindent
{\fontsize{9.3pt}{13.2pt}\selectfont
\textbf{\textcolor{BrandNavy}{Instituição Acadêmica Vinculada:}}\\
Universidade Federal de Goiás (UFG) — Instituto de Informática (INF)\\
Bacharelado em Engenharia de Software • Disciplina: Processos de Software\\[8pt]
\textbf{\textcolor{BrandNavy}{Diretoria Técnica e Autoria (Equipe Doomscoders):}}
\begin{listapontos}
  \itemponto{Carlos Eduardo Freire da Silva}{Arquitetura de Processos, Governança e Gestão de Portfólio}
  \itemponto{Gabriel Henrique Lourenço Dias}{Engenharia de Plataforma, Arquitetura Full-Stack e Automação CI/CD}
  \itemponto{Kayque Vicente Silva}{Engenharia de Requisitos, Qualidade de Software (QA) e Verificação}
  \itemponto{Yury Nogueira de Oliveira Silva}{Gerência de Configuração, Medição (GQM) e Melhoria Contínua}
\end{listapontos}
}

\vspace{6pt}
\noindent
{\fontsize{10pt}{13.5pt}\fontExtraBold\color{BrandNavy}Escopo Tecnológico do Produto de Referência (CineDoom):}\\[3pt]
{\fontsize{9.2pt}{13pt}\selectfont
O presente Manual regulamenta os processos técnicos, gerenciais e organizacionais aplicados ao desenvolvimento, operação e evolução da plataforma \textbf{CineDoom} — sistema SaaS e B2C de bilheteria digital, reserva concorrente de assentos em tempo real, catálogo de programação cinematográfica e vendas de bomboniere, estruturado em três macroentregas evolutivas (\textbf{v1.0 MVP Core}, \textbf{v2.0 Escala, Pagamentos PIX/Stripe e E-mail Transacional Resend} e \textbf{v3.0 Enterprise Multi-Tenant e Emissão Fiscal NFC-e}).}

\vspace{10pt}
\noindent
{\fontsize{10pt}{13.5pt}\fontExtraBold\color{BrandNavy}Histórico de Revisões do Manual:}\\[3pt]
\begin{center}
\small
\begin{tabularx}{\textwidth}{p{2.5cm} p{2.5cm} X}
\toprule
\textbf{Data} & \textbf{Versão} & \textbf{Descrição Sucinta da Alteração} \\
\midrule
10/09/2026 & Versão 1.0 & Criação do documento, identificação da equipe e elaboração da capa. \\
16/09/2026 & Versão 1.1 & Elaboração dos itens do Capítulo 1 (Apresentação do Software). \\
24/09/2026 & Versão 1.2 & Alinhamento dos Capítulos 1, 2 e 3 à norma ISO/IEC/IEEE 12207:2021. \\
24/09/2026 & Versão 1.3 & Revisão técnica: correção de contradições e justificativas de aquisição. \\
08/10/2026 & Versão 1.4 & Incremento e estruturação dos processos do Capítulo 4 (Gerenciamento). \\
17/10/2026 & Versão 1.5 & Reestruturação da capa e padronização do controle documental. \\
22/10/2026 & Versão 1.6 & Padronização formal da estrutura editorial dos Capítulos 1 a 5. \\
23/10/2026 & Versão 1.7 & Incorporação do Portal do Cliente, reserva ACID, e-mail e LGPD (ADR-012). \\
09/11/2026 & Versão 2.0 & Consolidação em 3 Macroentregas, matrizes Make-or-Buy, Pugh, EAP e Cap. 6 (MPS.BR). \\
\bottomrule
\end{tabularx}
\end{center}


\vspace{6pt}
\noindent
\begin{tcolorbox}[
  enhanced,
  colback=BrandTableAlt,
  colframe=BrandSecondary!60!BrandCrimson,
  boxrule=1.0pt,
  arc=6pt,
  left=12pt,
  right=12pt,
  top=9pt,
  bottom=9pt,
  fontupper=\fontsize{8.8pt}{12.4pt}\selectfont\color{BrandBodyText}
]
  {\fontsize{9.5pt}{12.5pt}\fontExtraBold\color{BrandNavy}Ficha Catalográfica e Controle Normativo do Documento}\\[5pt]
  \textbf{Título:} Manual de Processos de Software — CineSoft Solutions \& Plataforma CineDoom\\
  \textbf{Código de Controle de Configuração:} \codigoinline{CS-MPS-2026-V3}\\
  \textbf{Normas Técnicas de Referência:} ABNT NBR ISO/IEC/IEEE 12207:2021 • ISO/IEC 15504 (SPICE) • Guia Geral MPS de Software (MPS.BR / SOFTEX — Níveis G e F)\\
  \textbf{Stack de Engenharia Auditada:} Next.js 16 (App Router), React 19, TypeScript 5, Prisma ORM, PostgreSQL (Neon Serverless), Resend Transactional API, GitHub Actions e Vercel Cloud\\
  \textbf{Design System:} Cores Oficiais CineDoom (\#1A1A2E, \#0F3460, \#E94560, \#00D4FF)\\
  \textbf{Local e Ano de Publicação:} Goiânia — GO, Brasil • 2026
\end{tcolorbox}

\vfill
\noindent
\textcolor{BrandCrimson}{\rule{\textwidth}{0.8pt}}\\[3pt]
{\fontsize{8.2pt}{11pt}\color{BrandGrayMed}%
Template Editorial e Identidade Visual desenvolvidos em \LaTeX\ segundo a identidade de marca da CineSoft Solutions.}
\clearpage

% 3. LISTA DE ABREVIATURAS E SIGLAS
\tituloSecaoPrincipal{Lista de Abreviaturas e Siglas}

\noindent
{\fontsize{9pt}{13pt}\selectfont
\renewcommand{\arraystretch}{1.18}
\begin{tabularx}{\textwidth}{@{}>{\fontExtraBold\color{BrandNavy}}p{2.6cm} X@{}}
  ABNT      & Associação Brasileira de Normas Técnicas \\
  ADR       & \textit{Architecture Decision Record} (Decisão Arquitetural) \\
  API       & \textit{Application Programming Interface} \\
  B2B / B2C & \textit{Business-to-Business} / \textit{Business-to-Consumer} \\
  CI/CD     & \textit{Continuous Integration / Continuous Deployment} \\
  DoD / DoR & \textit{Definition of Done} / \textit{Definition of Ready} \\
  GPP       & Gerência de Portfólio de Projetos (Processo MPS.BR) \\
  GPR       & Gerência de Projetos (Processo MPS.BR — Nível G) \\
  GQM       & \textit{Goal-Question-Metric} (Paradigma de Medição) \\
  GQA       & Garantia da Qualidade (Processo MPS.BR — Nível F) \\
  GCO       & Gerência de Configuração (Processo MPS.BR — Nível F) \\
  IEC       & \textit{International Electrotechnical Commission} \\
  IEEE      & \textit{Institute of Electrical and Electronics Engineers} \\
  ISO       & \textit{International Organization for Standardization} \\
  JWT       & \textit{JSON Web Token} (Padrão de Autenticação RFC 7519) \\
  LGPD      & Lei Geral de Proteção de Dados Pessoais (Lei nº 13.709/2018) \\
  MED       & Medição (Processo MPS.BR — Nível F) \\
  MPS.BR    & Melhoria de Processo do Software Brasileiro (SOFTEX) \\
  MTTR      & \textit{Mean Time To Recovery} \\
  MVP       & \textit{Minimum Viable Product} (Produto Mínimo Viável) \\
  NFC-e     & Nota Fiscal de Consumidor Eletrônica (Modelo 65 — SEFAZ) \\
  ORM       & \textit{Object-Relational Mapping} (Prisma ORM) \\
  REQ       & Engenharia de Requisitos (Processo MPS.BR — Nível G) \\
  RNF       & Requisito Não Funcional \\
  SaaS      & \textit{Software as a Service} (Software como Serviço) \\
  SLA / SLO & \textit{Service Level Agreement} / \textit{Service Level Objective} \\
\end{tabularx}
}
\clearpage

% 4. LISTA DE FIGURAS E LISTA DE TABELAS/QUADROS
\tituloSecaoPrincipal{Lista de Figuras}
\makeatletter
\@starttoc{lof}
\makeatother

\vspace{22pt}
\tituloSecaoPrincipal{Lista de Tabelas e Quadros}
\makeatletter
\@starttoc{lot}
\makeatother
\clearpage

% 5. SUMÁRIO
\tituloSecaoPrincipal{Sumário}
\makeatletter
\@starttoc{toc}
\makeatother
\clearpage

% 6. APRESENTAÇÃO
\addMainEntryToTOC{Apresentação}
\tituloSecaoPrincipal{Apresentação}

A engenharia de sistemas digitais contemporâneos exige que organizações de tecnologia conciliem velocidade de entrega com previsibilidade operacional, segurança transacional e rastreabilidade normativa. No setor de entretenimento e exibição cinematográfica, onde picos de tráfego em estreias mundiais geram milhares de requisições simultâneas sobre um inventário finito de assentos físicos, a ausência de processos formais de ciclo de vida compromete diretamente a continuidade do negócio e a confiança do consumidor final.

Este documento constitui o \destaqueazul{Manual de Processos de Software da CineSoft Solutions}, elaborado pela equipe \termochave{Doomscoders} para institucionalizar a arquitetura de processos aplicada ao desenvolvimento, validação, entrega contínua e evolução da plataforma \destaqueazul{CineDoom}. Mais do que um referencial descritivo, este manual consolida práticas reais de engenharia de software alinhadas simultaneamente ao padrão internacional \termochave{ABNT NBR ISO/IEC/IEEE 12207:2021} e ao modelo brasileiro de maturidade \termochave{MPS.BR (SOFTEX)} nos níveis \textbf{G (Parcialmente Gerenciado)} e \textbf{F (Gerenciado)}.

Para garantir aderência à realidade operacional de uma empresa de produto digital, a estratégia de engenharia documentada neste manual estrutura a evolução do ecossistema CineDoom em \textbf{três macroentregas progressivas}:
\begin{listapontos}
  \itemponto{Macroentrega 1 (CineDoom v1.0 — MVP Core Operacional)}{Compreende o núcleo transacional construído e implantado em nuvem, abrangendo catálogo de filmes e sessões, máquina de estados de reserva de assentos com bloqueio temporário (\codigoinline{SeatHold}), compra de ingressos e bomboniere, autenticação JWT e painel administrativo.}
  \itemponto{Macroentrega 2 (CineDoom v2.0 — Escala, Pagamentos Reais e Mensageria)}{Expande a infraestrutura para alta concorrência, integrando gateway de pagamento real (PIX dinâmico com webhook e Stripe), serviço de e-mail transacional \textbf{Resend} para envio automatizado de vouchers QR Code e recuperação de senha, além de banco PostgreSQL gerenciado com alta disponibilidade.}
  \itemponto{Macroentrega 3 (CineDoom v3.0 — Plataforma SaaS Multi-Tenant e Fiscal)}{Consolida o modelo de negócios B2B2C para múltiplas redes exibidoras, incorporando emissão fiscal automatizada (NFC-e junto à SEFAZ), motor de precificação dinâmica, programa de fidelidade e observabilidade distribuída com OpenTelemetry.}
\end{listapontos}

Ao longo das unidades deste volume, são apresentados os fundamentos organizacionais, o mapeamento normativo, os fluxos detalhados de cada processo (entradas, atividades, papéis, artefatos e critérios de verificação), bem como os mecanismos de medição GQM e melhoria contínua que sustentam a excelência técnica da CineSoft Solutions.
\clearpage
